\documentclass[12pt, a4paper]{article}

\usepackage[T1]{fontenc}
\usepackage[utf8]{inputenc}
\usepackage[british]{babel}
\usepackage{mathtools}
\usepackage{amsthm}
\usepackage{libertine}
\usepackage[libertine]{newtxmath}
\usepackage[mathscr]{euscript}
\usepackage{listings}
\usepackage{enumitem}
\usepackage{tikz-cd}
\usetikzlibrary{decorations.markings}
\usepackage{float}
\usepackage{geometry}
 \geometry{
 a4paper,
 total={170mm,257mm},
 left=20mm,
 top=30mm,
 }
\usepackage{fancyhdr}
\usepackage{hyperref}
\urlstyle{same}
\usepackage[noabbrev]{cleveref}
\usepackage{xcolor}
\usepackage{tcolorbox}
\tcbuselibrary{listings}

% Colors to match VS Code Lean 4 Light Theme
\definecolor{keywordcolor}{rgb}{0.0, 0.0, 1.0}   % blue
\definecolor{tacticcolor}{rgb}{0.0, 0.0, 1.0}    % blue
\definecolor{commentcolor}{rgb}{0.0, 0.5, 0.0}   % green
\definecolor{symbolcolor}{rgb}{0.0, 0.0, 0.0}    % black
\definecolor{sortcolor}{rgb}{0.0, 0.0, 1.0}      % blue
\definecolor{attributecolor}{rgb}{0.0, 0.0, 0.0} % black
% Light gray background color
\definecolor{codebg}{rgb}{0.97, 0.97, 0.97}

\def\lstlanguagefiles{lstlean.tex}

\newtcblisting{leancode}{
  listing only,
  listing options={
    language=lean,
    emph={sorry},
    emphstyle=\color{red},
    basicstyle=\ttfamily\footnotesize
  },
  colback=codebg,
  colframe=codebg,
  boxrule=0pt,
  arc=3mm,
  auto outer arc,
  top=10pt,
  bottom=10pt,
  left=5pt,
  right=5pt,
  width=\textwidth,
}

\theoremstyle{definition}
\newtheorem{innercustomthm}{Exercise}
\newenvironment{xca}[1] {\renewcommand\theinnercustomthm{#1}\innercustomthm}
  {\endinnercustomthm}

\pagestyle{fancy}
\fancyhf{}
\rhead{Lean Algebra Exercises{\renewcommand{\thefootnote}{\fnsymbol{footnote}}\footnotemark[2]}}
\chead{Normal Subgroup}
\lhead{Pedro N\'{u}\~{n}ez}
% \rfoot{Page \thepage}

\title{Lean Algebra Exercise 2.1}
\author{Pedro N\'{u}\~{n}ez}

\setcounter{tocdepth}{1}
\sloppy
\makeatletter
\hypersetup{
  pdfauthor={\@author},
  pdftitle={\@title},
  colorlinks,
  linkcolor=[rgb]{0.2,0.2,0.6},
  citecolor=[rgb]{0.2,0.2,0.6},
  urlcolor=[rgb]{0.2,0.2,0.6}}
\makeatother

\begin{document}

\begin{xca}{2.1}
  Let $H \subset G$ be a subgroup.
  Show that the following assertions are equivalent:
  \begin{enumerate}[label=(\alph*)]
    \item For all $g \in G$, $Hg = gH$.
    \item For all $g \in G$, $g^{-1}Hg = H$.
    \item For all $g \in G$, $g^{-1}Hg \subset H$.
    \item For all $g_{1}, g_{2} \in G$, $g_{1}g_{2} \in H$ if and only if $g_{2}g_{1} \in H$.
  \end{enumerate}

  You can try to \href{https://live.lean-lang.org/#codez=JYWwDg9gTgLgBAWQIYwBYBtgCMB0BBdAcwFMsokcBxKCAVzBwGVatCb6cAhJAZ2AGMAUKEixEKDNhwAVJPxgCZAMTwBRQYIBuSKMCRZ0xOAG9KcAFxxpATzDEAvnADa1OmDiUAunAAUAOQs4ZlZ2d0oASg0AE2IAMzgAO2gQJHQAfSQLQTg4AAUad3MAXjhAACI4QkDKABoyuFRAv1rAYCJ6gHJG2sBEwgaAKgq4Esr%2B1DbouMTk1LSsLJz8iEKS8srLGt9yhssmuFbRzorAbwIATrgRs4GS0fC4QHIibI36g72O7e6GoZOL%2Ff7CSMEYvEklAUul%2BHM8gULMsBmtaptnu0DoQvudfoN2uMgVN0lEIQslnVCABGCoAJiqtR8e2RpN%2BFKuN0AKYS%2BGlvckXEkY1D%2FQwgFJwJBpESYYg8GYQ8r%2BQLBNhuDzhWrA0EZOABQBJhJMQdNZgFinAsNYHsAEjAaDw1YKeDxaOAYA9aAlYhB0HjldNMigrTa7Q8eNAoEbBGhiNBiCA4DBYkhiGl3ekbSF5eYHgAZYA8GDKNTOeOqnZ53VKnFpcEFkt4vzefWGh5RmNpVBITRGUmain6oUi4BiiV%2BOvR2NNltwCmagDM0INRpyOX9UEDA4bw6Mk81ABYp7XZ3B54ucvWh82jJvNaSazPZ3vL4e0rETRnUIIgA}{verify your solution to this exercise with Lean}.
\end{xca}

\subsection*{Hints for Lean}

\begin{leancode}
import Mathlib.Algebra.Group.Subgroup.Basic
import Mathlib.Tactic.TFAE

variable {G : Type} [Group G] (N : Subgroup G)

def normal_a :
  Prop := ∀ g : G, ∀ h : N, ∃ h' : N, ↑h * g = g * h'

def normal_b :
  Prop := ∀ g : G, (∀ h : N, ∃ h' : N, g⁻¹ * h * g = h') ∧
  (∀ h : N, ∃ h' : N, ↑h = g⁻¹ * h' * g)

def normal_c :
  Prop := ∀ g : G, ∀ h : N, ∃ h' : N, g⁻¹ * h * g = h'

def normal_d :
  Prop := ∀ g1 g2 : G, (∃ h : N, g1 * g2 = h) ↔ (∃ h : N, g2 * g1 = h)

lemma a_implies_b :
  ∀ (N : Subgroup G), normal_a N → normal_b N := by
  intros N assumpt
  unfold normal_a at assumpt
  sorry

theorem tfae_normal_subgroup :
  List.TFAE [normal_a N, normal_b N, normal_c N, normal_d N] := by
  tfae_have 1 → 2 := a_implies_b N
  tfae_have 2 → 3 := by
    sorry
  tfae_have 3 → 4 := by
    sorry
  tfae_have 4 → 1 := by
    sorry
  tfae_finish
\end{leancode}

You may find the \href{https://leanprover-community.github.io/mathlib4_docs/Mathlib/Algebra/Group/Subgroup/Basic.html}{subgroup documentation} on mathlib helpful.
The vertical arrow $\uparrow$ {\ttfamily h} also shows up \href{https://github.com/leanprover-community/tutorials4/blob/master/Tutorials/Solutions/09LimitsFinal.lean}{somewhere in the lean 4 tutorials project}.
A \href{https://github.com/pedro-nunez/lean-algebra-2021/blob/master/normal-subgroup/solution.lean}{solution} is available on the repository
\begin{center}
  \href{https://github.com/pedro-nunez/lean-algebra-2021}{https://github.com/pedro-nunez/lean-algebra-2021}.
\end{center}

{\renewcommand{\thefootnote}{\fnsymbol{footnote}}\footnotetext[2]{Exercises designed for the lecture \href{https://home.mathematik.uni-freiburg.de/arithgeom/lehre/ws21/azt}{\emph{Algebra and Number Theory}}, for which I was a teaching assistant.
This was an introductory lecture on abstract algebra, polynomial rings and Galois theory, taught by Annette Huber-Klawitter at the University of Freiburg during the Winter Semester 2021/2022.}}

\end{document}
